\documentclass[a4paper,11pt]{article}
\usepackage[utf8]{inputenc}
\usepackage[T1]{fontenc}
\usepackage[french]{babel}
\usepackage{geometry}
\geometry{left=2cm,right=2cm,top=1.5cm,bottom=2cm}
\usepackage{amsmath,amssymb}
\usepackage{enumitem}
\usepackage{array}
\usepackage{booktabs}
\usepackage{xcolor}
\usepackage{tcolorbox}
\usepackage{fancyhdr}
\usepackage{tikz}
\usepackage{lastpage}

% Couleurs
\definecolor{bleupro}{HTML}{0969da}
\definecolor{orange}{HTML}{d97706}
\definecolor{vertpro}{HTML}{2f855a}
\definecolor{rouge}{HTML}{c53030}
\definecolor{gris}{HTML}{718096}

% En-tête
\pagestyle{fancy}
\fancyhf{}
\lhead{\small CCF Mathématiques — \textbf{CORRECTION}}
\rhead{\small Terminale Bac Pro ICCER}
\cfoot{\thepage/\pageref{LastPage}}
\renewcommand{\headrulewidth}{0.5pt}

% Boîtes
\tcbuselibrary{skins,breakable}

\newtcolorbox{corrbox}{
  colback=green!3,
  colframe=vertpro,
  title={\textbf{Correction}},
  fonttitle=\bfseries,
  breakable
}

\newtcolorbox{contextebox}{
  colback=blue!3,
  colframe=bleupro,
  leftrule=4pt,
  rightrule=0pt,
  toprule=0pt,
  bottomrule=0pt,
  boxsep=4pt,
  left=8pt
}

\newcommand{\pts}[1]{\hfill\textcolor{gris}{\small\textbf{#1}}}
\newcommand{\comp}[1]{\colorbox{blue!10}{\scriptsize\textbf{#1}}}

\setlength{\parindent}{0pt}
\setlength{\parskip}{6pt}

\begin{document}

\begin{center}
{\LARGE\bfseries\color{rouge} CORRECTION — CCF Mathématiques}\\[8pt]
{\large Terminale Bac Pro — ICCER}\\[8pt]
{\Large\bfseries Probabilités \& Fonctions polynômes de degré 3}\\[12pt]
\textcolor{rouge}{\rule{12cm}{1pt}}
\end{center}

\vspace{10pt}

% ================================================================
\section*{Exercice 1 — Contrôle qualité en génie climatique \pts{5 points}}
% ================================================================

\begin{contextebox}
\textbf{Données :} $P(E) = 0{,}10$ \quad $P(\overline{E}) = 0{,}90$ \quad $P_E(F) = 0{,}25$ \quad $P_{\overline{E}}(F) = 0{,}04$

\smallskip
$E$ = ``défaut d'étanchéité'' \quad $F$ = ``défaut de filetage''
\end{contextebox}

\vspace{6pt}

% --- Q1 ---
\textbf{Q1} \comp{APP} — Traduire les données en probabilités. \pts{1 pt}

\begin{corrbox}
D'après l'énoncé :
\begin{itemize}[nosep]
  \item $P(E) = 0{,}10$ \quad (10\% des raccords ont un défaut d'étanchéité)
  \item $P(\overline{E}) = 1 - P(E) = 1 - 0{,}10 = \boxed{0{,}90}$
  \item $P_E(F) = 0{,}25$ \quad (probabilité de défaut de filetage sachant défaut d'étanchéité)
  \item $P_{\overline{E}}(F) = 0{,}04$ \quad (probabilité de défaut de filetage sachant pas de défaut d'étanchéité)
\end{itemize}

\textit{Barème : 0,25 pt par valeur correcte.}
\end{corrbox}

% --- Q2 ---
\textbf{Q2} \comp{REA} — Compléter l'arbre de probabilités. \pts{1 pt}

\begin{corrbox}
\begin{center}
\begin{tikzpicture}[scale=1, every node/.style={font=\small}]
  \node[circle, fill=black, inner sep=1.5pt] (R) at (0,0) {};
  % Niveau 1
  \node (E) at (3,1.5) {$E$};
  \node (Eb) at (3,-1.5) {$\overline{E}$};
  \draw[-] (R) -- (E) node[midway, above, sloped] {$0{,}10$};
  \draw[-] (R) -- (Eb) node[midway, below, sloped] {$0{,}90$};
  % Niveau 2 depuis E
  \node (EF) at (6.5,2.3) {$F$};
  \node (EFb) at (6.5,0.7) {$\overline{F}$};
  \draw[-] (E) -- (EF) node[midway, above, sloped] {$0{,}25$};
  \draw[-] (E) -- (EFb) node[midway, below, sloped] {$0{,}75$};
  % Niveau 2 depuis Eb
  \node (EbF) at (6.5,-0.7) {$F$};
  \node (EbFb) at (6.5,-2.3) {$\overline{F}$};
  \draw[-] (Eb) -- (EbF) node[midway, above, sloped] {$0{,}04$};
  \draw[-] (Eb) -- (EbFb) node[midway, below, sloped] {$0{,}96$};
\end{tikzpicture}
\end{center}

\textbf{Vérifications :}
\begin{itemize}[nosep]
  \item Nœud racine : $0{,}10 + 0{,}90 = 1$ \checkmark
  \item Nœud $E$ : $0{,}25 + 0{,}75 = 1$ \checkmark \quad ($P_E(\overline{F}) = 1 - 0{,}25 = 0{,}75$)
  \item Nœud $\overline{E}$ : $0{,}04 + 0{,}96 = 1$ \checkmark \quad ($P_{\overline{E}}(\overline{F}) = 1 - 0{,}04 = 0{,}96$)
\end{itemize}

\textit{Barème : 0,5 pt pour les branches niveau 1 + 0,5 pt pour les branches niveau 2.}
\end{corrbox}

% --- Q3 ---
\textbf{Q3} \comp{REA} — Calculer $P(E \cap F)$ et $P(\overline{E} \cap F)$. \pts{1 pt}

\begin{corrbox}
On multiplie les probabilités le long de chaque chemin :

\textbf{Chemin $E$ puis $F$ :}
\[
P(E \cap F) = P(E) \times P_E(F) = 0{,}10 \times 0{,}25 = \boxed{0{,}025}
\]

\textbf{Chemin $\overline{E}$ puis $F$ :}
\[
P(\overline{E} \cap F) = P(\overline{E}) \times P_{\overline{E}}(F) = 0{,}90 \times 0{,}04 = \boxed{0{,}036}
\]

\textit{Barème : 0,5 pt par calcul correct.}
\end{corrbox}

% --- Q4 ---
\textbf{Q4} \comp{ANA} — Calculer $P(F)$ par la formule des probabilités totales. \pts{1 pt}

\begin{corrbox}
La formule des probabilités totales donne :
\[
P(F) = P(E \cap F) + P(\overline{E} \cap F) = 0{,}025 + 0{,}036 = \boxed{0{,}061}
\]

Soit \textbf{6,1\%} des raccords ont un défaut de filetage.

\smallskip
\textit{Explication :} l'événement $F$ peut se produire par deux chemins dans l'arbre (via $E$ ou via $\overline{E}$). La probabilité totale est la somme des probabilités de ces deux chemins.

\smallskip
\textit{Barème : 0,5 pt pour la formule posée + 0,5 pt pour le résultat.}
\end{corrbox}

% --- Q5 ---
\textbf{Q5} \comp{ANA} — Application et conclusion. \pts{1 pt}

\begin{corrbox}
Sur 500 raccords :
\[
\text{Nombre attendu de raccords avec défaut de filetage} = 500 \times 0{,}061 = \boxed{30{,}5 \approx 31 \text{ raccords}}
\]

Taux de défaut de filetage : $\frac{30{,}5}{500} = 6{,}1\%$.

Le responsable qualité considère le taux acceptable si $< 5\%$. Or $6{,}1\% > 5\%$.

$\Rightarrow$ \textbf{Non}, le taux n'est pas acceptable. Il faudra améliorer le processus de fabrication pour réduire les défauts, en particulier les défauts de filetage sur les raccords ayant déjà un défaut d'étanchéité ($P_E(F) = 25\%$ est très élevé).

\smallskip
\textit{Barème : 0,5 pt pour le calcul + 0,5 pt pour la conclusion argumentée.}
\end{corrbox}

\newpage

% ================================================================
\section*{Exercice 2 — Optimisation du volume d'un bac de rangement \pts{5 points}}
% ================================================================

\begin{contextebox}
\textbf{Données :} Plaque carrée 30 cm $\times$ 30 cm. Découpe de côté $x$ aux coins.\\
$V(x) = x(30-2x)^2 = 4x^3 - 120x^2 + 900x$ pour $x \in \;]0\,;\,15[$
\end{contextebox}

\vspace{6pt}

% --- Q1 ---
\textbf{Q1} \comp{REA} — Tracer la courbe à la calculatrice et la reproduire sur le quadrillage. \pts{1 pt}

\begin{corrbox}
Sur la calculatrice NumWorks (ou GeoGebra), on saisit $f(x) = 4x^3 - 120x^2 + 900x$ et on trace sur $[0\,;\,15]$.

\begin{center}
\begin{tikzpicture}[scale=0.48]
  % Quadrillage
  \draw[gray!25, very thin] (0,0) grid[step=1] (16,14);
  \draw[gray!50, thin] (0,0) grid[step=5] (16,14);
  % Axes
  \draw[->, thick] (0,0) -- (16.5,0) node[right] {$x$ (cm)};
  \draw[->, thick] (0,0) -- (0,14.5) node[above] {$V$ (cm$^3$)};
  % Graduations x
  \foreach \x in {5,10,15} {
    \node[below] at (\x,-0.3) {\small\x};
  }
  % Graduations y (1 unité graphique = 250 cm³)
  \node[left] at (-0.3,2) {\small 500};
  \node[left] at (-0.3,4) {\small 1000};
  \node[left] at (-0.3,6) {\small 1500};
  \node[left] at (-0.3,8) {\small 2000};
  \node[left] at (-0.3,10) {\small 2500};
  \node[left] at (-0.3,12) {\small 3000};
  \node[below left] at (0,0) {\small 0};
  % Courbe V(x) = 4x³ - 120x² + 900x (échelle y : /250)
  \begin{scope}
    \clip (0,0) rectangle (15,14);
    \draw[thick, bleupro, smooth] plot[domain=0:15, samples=50]
      (\x, {(4*\x*\x*\x - 120*\x*\x + 900*\x)/250});
  \end{scope}
  % Maximum en x=5, V(5)=2000
  \filldraw[rouge] (5,8) circle (3pt);
  \node[above right, rouge, font=\small\bfseries] at (5.2,8) {max $(5\,;\,2\,000)$};
  % Pointillés de lecture
  \draw[dashed, rouge, thin] (5,0) -- (5,8);
  \draw[dashed, rouge, thin] (0,8) -- (5,8);
\end{tikzpicture}
\end{center}

L'élève doit reproduire cette allure sur le quadrillage fourni, avec le maximum repéré.

\smallskip
\textit{Barème : 0,5 pt pour la forme générale correcte + 0,5 pt pour le maximum placé.}
\end{corrbox}

% --- Q2 ---
\textbf{Q2} \comp{APP} — Lecture graphique du maximum. \pts{1 pt}

\begin{corrbox}
Par lecture graphique sur la calculatrice :

Le volume semble maximal pour $x \approx \boxed{5 \text{ cm}}$.

La valeur approchée du maximum est $V \approx \boxed{2\,000 \text{ cm}^3}$.

\smallskip
\textit{Barème : 0,5 pt pour la valeur de $x$ (tolérance $\pm 0{,}5$) + 0,5 pt pour la valeur de $V$ (tolérance $\pm 200$).}
\end{corrbox}

% --- Q3 ---
\textbf{Q3} \comp{REA} — Dérivée et résolution. \pts{1 pt}

\begin{corrbox}
\textbf{Dérivée :}
\[
V'(x) = 3 \times 4x^2 - 2 \times 120x + 900 = 12x^2 - 240x + 900
\]

\textbf{Vérification de la factorisation :}
\[
12(x-5)(x-15) = 12(x^2 - 20x + 75) = 12x^2 - 240x + 900 \quad \checkmark
\]

\textbf{Résolution} $V'(x) = 0$ :
\[
x - 5 = 0 \quad \Rightarrow \quad \boxed{x = 5}
\]
\[
x - 15 = 0 \quad \Rightarrow \quad x = 15 \quad \text{(borne de l'intervalle, exclue)}
\]

Seule solution dans $]0\,;\,15[$ : $\boxed{x = 5 \text{ cm}}$

\smallskip
\textit{Barème : 0,5 pt pour la dérivée + 0,5 pt pour la résolution.}
\end{corrbox}

% --- Q4 ---
\textbf{Q4} \comp{ANA} — Compléter le tableau de variations. Calculer le maximum exact. \pts{1 pt}

\begin{corrbox}
\textbf{Tableau de variations complété :}

\begin{center}
\renewcommand{\arraystretch}{1.6}
\begin{tabular}{|>{\centering\arraybackslash}m{2.5cm}|>{\centering\arraybackslash}m{1.2cm}|>{\centering\arraybackslash}m{1.8cm}|>{\centering\arraybackslash}m{1.2cm}|>{\centering\arraybackslash}m{1.8cm}|>{\centering\arraybackslash}m{1.2cm}|}
\hline
$x$ & $0$ & & $5$ & & $15$ \\
\hline
Signe de $(x - 5)$ & $-$ & $\boldsymbol{-}$ & $0$ & $\boldsymbol{+}$ & $+$ \\
\hline
Signe de $(x - 15)$ & $-$ & $\boldsymbol{-}$ & $-$ & $\boldsymbol{-}$ & $0$ \\
\hline
Signe de $V'(x)$ & & $\boldsymbol{+}$ & $0$ & $\boldsymbol{-}$ & \\
\hline
& $0$ & & $\boldsymbol{2\,000}$ & & $0$ \\[-2pt]
Variations de $V$ & & $\nearrow$ & & $\searrow$ & \\
\hline
\end{tabular}
\renewcommand{\arraystretch}{1}
\end{center}

\textbf{Règle des signes :} $V'(x) = 12 \times (x-5) \times (x-15)$. Le facteur 12 est positif, donc le signe de $V'$ est celui du produit $(x-5)(x-15)$ :
\begin{itemize}[nosep]
  \item $(-) \times (-) = +$ sur $]0\,;\,5[$ → $V$ croissante
  \item $(+) \times (-) = -$ sur $]5\,;\,15[$ → $V$ décroissante
\end{itemize}

\textbf{Calcul du maximum :}
\[
V(5) = 4 \times 125 - 120 \times 25 + 900 \times 5 = 500 - 3\,000 + 4\,500 = \boxed{2\,000 \text{ cm}^3}
\]

\textbf{Confirmation :} la lecture graphique de Q2 donnait $V \approx 2\,000$ cm$^3$ pour $x \approx 5$. Le calcul confirme exactement \checkmark

\smallskip
\textit{Barème : 0,5 pt pour le tableau de variations complété + 0,5 pt pour le calcul du maximum.}
\end{corrbox}

% --- Q5 ---
\textbf{Q5} \comp{COM} — Conclusion et dimensions. \pts{1 pt}

\begin{corrbox}
Le volume est maximal pour $\boxed{x = 5 \text{ cm}}$.

\textbf{Dimensions du bac :}
\begin{itemize}[nosep]
  \item Hauteur : $x = 5$ cm
  \item Longueur = Largeur : $30 - 2 \times 5 = 20$ cm
\end{itemize}

Le bac est un parallélépipède rectangle de \textbf{20 cm $\times$ 20 cm $\times$ 5 cm}.

\textbf{Volume maximal :} $V_{\max} = 20 \times 20 \times 5 = \boxed{2\,000 \text{ cm}^3}$

\textbf{Contrainte :} L'installateur a besoin d'un bac d'au moins 1\,500 cm$^3$. Or $2\,000 > 1\,500$.

$\Rightarrow$ \textbf{Oui, la contrainte est respectée.} Le bac optimal de 2\,000 cm$^3$ dépasse largement le minimum requis.

\smallskip
\textit{Barème : 0,25 pt dimensions + 0,25 pt volume max + 0,5 pt conclusion sur la contrainte.}
\end{corrbox}

\vfill

\begin{center}
\textcolor{rouge}{\rule{12cm}{1pt}}\\[8pt]
\begin{tabular}{|c|c|c|c|}
\hline
\textbf{Exercice} & \textbf{Thème} & \textbf{Questions} & \textbf{Points} \\
\hline
1 & Probabilités (raccords cuivre, arbre, probas totales) & Q1 à Q5 & \textbf{/5} \\
\hline
2 & Polynôme degré 3 (dérivée, variations, optimisation) & Q1 à Q5 & \textbf{/5} \\
\hline
\multicolumn{3}{|r|}{\textbf{Total}} & \textbf{/10} \\
\hline
\end{tabular}
\end{center}

\end{document}
